<<<P000922>[2].
Иначе
<<<P000923>{(при отсутствии циклов зависимости)},
<<<p000010> быть наименьшим числом, таким, что <<<P000011> встречается в <<<P000012> Для некоторых <<P000013>
и <<<P000014> для всех <<<P000015> в <<<P000016>
<<<P000924>{%>
(то есть, граница <<<P000017>) не содержит переменных типа,
<<<p000018> возникает в связи с некоторой другой переменной типа )%
}
Для всех остальных (P000019).
<<<P000020> Получено из <<<P000021>
Заменить <<P000022> Для каждого случая <<P000023>
Находится в позиции <<<P000024> который не является контравариантным,
и замена <<<P000584> для каждого случая <<<P000025>
который находится в противоположном положении в <<<P000026>.

<<<P000925>[3].
Иначе
<<<P000926>{(когда вообще нет зависимостей)},
Окончание с результатом <<<P000585>.
<<P000485>

<<<P000927>{%>
Этот процесс всегда заканчивается, так как общее количество
Возникновение переменных типа из <<<P000028> в
Текущие границы строго уменьшаются с каждым шагом, и мы заканчиваем.
когда это число достигает нуля.%
?

<<<P000928>{%>
Может показаться несколько произвольным рассматривать неиспользованные и инвариантные параметры.
аналогично ковариантным параметрам,
В частности, из-за неисправности инвариантных параметров
Чтобы оправдать ожидания, что
Необработанный тип обозначает супертип всех выразительных типов с регулярными границами.

Мы могли бы легко сделать все возможное, чтобы связать ошибку.
когда применяется к типу, где инвариантность возникает в любом месте
во время выполнения алгоритма.
Тем не менее, есть ряд случаев, когда этот выбор производит полезный тип.
И мы решили, что не стоит запрещать такие случаи.%
?

<<P000000>

<<<P000929>{%>
Например, значение <<<P000586> может иметь динамический тип
<<P000587>.
Тем не менее, аргументы типа должны быть тщательно подобраны.
или результат не будет подтипом <<<P000588>.
Например, <<P000589> не может иметь динамический тип
<<P000930>>>{B<int, Inv<int>{}>}
<<<P000590>> не является подтипом <<<P000591>.%
?

<<<P000837>%
Необработанный тип <<<P000029> является ошибкой времени компиляции, если инстанциация связана с <<<P000030>
Устанавливает тип, который не является хорошо ограниченным
(P000514).

<<<P000931><%
Такая ошибка может произойти, как показано на следующем примере:
?

<<<P000001>

<<<P000932>{%>
С этими заявлениями,
<<<P000592> В качестве аннотации типа используется ошибка времени компиляции:
Алгоритм дает <<P000933> F<C\DYNAMIC{}>
И это не регулярный и не сверхограниченный тип.
%
Полученный тип может быть определен явно как
\code{C\DYNAMIC{} <<P000937>>>(C\DYNAMIC{})}.
Такой тип существует,
Мы просто не можем выразить это, передав аргумент типа <<<P000593>,
Таким образом, мы делаем ошибку, а не допускаем ее неявно.
?

<<<P000939>{%>
Основная причина, почему имеет смысл сделать такой сырой тип ошибкой
Дело в том, что не существует субтипа отношений.
между соответствующими параметризованными типами.%
?
<<<P000940>{%>
Например, <<<P000594> и <<<P000595> не связаны между собой,
Даже когда \SubtypeNE{\code{T1}}{\code{T2}> или наоборот.
На самом деле, нет типа <<P000598> какой бы то ни было
переменная с заявленным типом <<P000599>
может быть присвоен переменной типа
\code{C\DYNAMIC{} \FUNCTION(C\DYNAMIC{})}.
%
Таким образом, сырой <<<P000600>, если это разрешено,
Это не было бы супертипом P000601. для всех возможных <<<P000602>>,
Это будет тип, который не связан с <<<P000603>
<<<P000946>{каждый сингл} <<P000604> удовлетворяет пределу <<<P000605>.
Это настолько бесполезно, что мы сделали ошибку.
?

<<<P000838>%
Когда инстанциация в связь применяется к типу, она протекает рекурсивно:Для параметризированного типа <<<P000947><<<P000031><<<<P000948> T}{1}{k}>
применяется к <<<P000949> T}{1}{k}
Для функционального типа
<<P000950> T_0

<<P000951>
и тип функции
<<P000952> T_0
применяется к <<<P000953> T}{0}{n+k}.

<<<P000954>{%>
Это означает, что установление связи не влияет на
Тип, который не содержит сырых типов.
И наоборот, инстанциация к связанным действиям на типы, которые являются синтаксическими субтерминами.
когда они глубоко вложены.%
?

<<<P000839>%
<<P000955> Обоснование связи с общей функцией <<<P000032>{%}
Общая функция! моментализация для связывания
использует алгоритм, описанный выше,
Взять формальные параметры <<P000956> X}{1}{k) из декларации <<P000033>,
с границами <<<P000957>{B}{1}{k} и,
для каждого <<P000034>,
Позволяет <<<P000035> обозначает результат инстанциации, связанный с <<<P000036>,
и позволить <<<P000037>\DYNAMIC{} когда граница <<<P000038> опущена.

<<<P000840>%
Пусть <<<P000039> будет родовой декларацией функции.
Если инстанция связана с <<<P000040> доходность
Список типов аргументов <<P000959> T}{1}{k) такой,
для некоторых <<<P000041>,
<<P000042> является или содержит тип, который не является хорошо ограниченным,
или если <<P000960> T}{1}{k} не удовлетворяет границам
по формальным параметрам типа <<<P000043>,
Мы говорим, что
<<<P000961><<<P000044> не имеет аргументов типа по умолчанию {%}
Общая функция! не имеет аргументов типа по умолчанию.


<<<P000962>{Метадата}
<<P000908>

<<<P000841>%
Dart поддерживает метаданные, которые используются для прикрепления
Пользователь определяет аннотации к структуре программы.

<<P000457>
<metadata> ::= ('@' <metadatum>)*

<metadatum> ::= <<<P000963><
<identifier> | <qualifiedName> | <constructorDesignation> <аргументы>
<<P000486>

<<<P000842>%
Метаданные состоят из серии аннотаций,
Каждый из них начинается с персонажа. <<<P000964>{@}
За ним следует постоянное выражение <<P000045> извлекаемый из
<<<P000965>{metadatum}.
Ошибка времени компиляции <<<P000046> не является одним из следующих:

<<P000458>
<<P000966> Ссылка на постоянную переменную.
<<P000967> Призыв к постоянному конструктору.
<<P000487>

<<<P000968>{%>
Выражение <<P000047> происходит в постоянном контексте
(<<P000515>),
Это означает, что <<P000969> Модификаторы не должны быть указаны явно.%
?

<<<P000843>%
Метаданные, которые являются первым элементом
Правая сторона грамматического правила
связано с абстрактным синтаксическим деревом для нетерминального
на левой стороне грамматического правила
<<<P000970>{(то есть он связан с его родителем)}.
В противном случае метаданные связаны с абстрактным деревом синтаксиса.
Конструкция программы <<<P000048> что сразу следует за метаданными
Правило грамматики.

<<<P000971>{%>
Эти правила призваны обеспечить минимальный уровень согласованности.
в ассоциации, которая связывает каждый метадатум с конструкцией программы.
Структура абстрактного дерева синтаксиса специфична для реализации.
И даже не гарантируется, что инструмент будет использовать что-либо.
Это называется абстрактным синтаксическим деревом.
В этом случае выражение "абстрактное синтаксическое дерево" следует толковать как
Субъекты представления программ, которые фактически используются.

Это означает, что тонкие детали связи между метаданными
Абстрактный узел дерева синтаксиса также является специфическим для реализации.
В частности, реализация может выбрать ассоциацию данного метадатума.
с более чем одним абстрактным узлом дерева синтаксиса.%
?

<<<P000844>%
Метаданные могут быть получены во время выполнения с помощью отражающего вызова,
Предоставляется аннотированная конструкция программы <<<P000049> доступен через отражение.

<<<P000972>{%>
Метаданные также могут быть получены статически.
Анализ программы и оценка констант с помощью подходящего интерпретатора.
На самом деле, многие, если не большинство применений метаданных полностью статичны.В этом случае связывание каждого метадатума с абстрактным синтаксическим узлом дерева
определяется данным инструментом статического анализа,
которые, разумеется, не подпадают под какие-либо ограничения настоящего документа.
Конечно, будет полезно стремиться к согласованности всех инструментов.
относительно связывания метаданных с абстрактными узлами синтаксиса.%
?

<<<P000973><%
Важно, чтобы накладные расходы по времени выполнения не были понесены
Введение метаданных, которые фактически не используются.
Поскольку метаданные содержат только константы,
Время, в которое он вычисляется, не имеет значения.
Таким образом, реализация может пропустить метаданные.
при обычном разборе и исполнении,
и оценивать его лениво.%
?

<<<P000974>{%>
Метаданные можно связать с конструктами
Это может быть недоступно через рефлексию.
Например, локальные переменные
(хотя можно предположить, что в будущем,
Более богатые отражающие библиотеки также могут обеспечить доступ к ним.
Это не так бесполезно, как может показаться.
Как отмечалось выше, данные могут быть получены статически.
Если исходный код доступен.%
?

<<<P000845>%
Метаданные могут появляться перед библиотекой, заголовком части, классом,
типдеф, параметр типа, конструктор, завод, функция,
параметр или переменная декларация,
перед импортом, экспортом или частичной директивой.

<<<P000846>%
Постоянное выражение, данное в аннотации, проверяется и оценивается по типу.
в объеме, окружающем аннотируемое заявление.


<<P000975> Выражения
<<P000909>

<<<P000847>%
<<<P000570> Это фрагмент кода Дарта.
Это можно оценить во время выполнения.

<<<P000848>%
Каждое выражение имеет ассоциированный статический тип (<<P000516>).
может иметь ассоциированный статический тип контекста
% % TODO(Eernst): Это не определено до тех пор, пока CL 92782 не приземлится.
% (<<P000517>),
Это может повлиять на статический тип и оценку выражения.
Каждый объект имеет соответствующий динамический тип (P000518).

% % TODO(eernst): <expression> следует вывести <functionExpression> Кроме того,
% от того, что в настоящее время получено из <первичного> Побуждает к подлинному
%% и серьезный источник двусмысленности, что затрудняет разбор Дарта
% с использованием чего-либо, кроме модифицированного рекурсивного парсера спуска.

<<P000459>
<expression> ::= <assignableExpression> <assignmentOperator> <expression>
<<P000976> <условное выражение>
<<P000977> <cascade>
<<P000978> <rowExpression>

<expressionWithoutCascade> ::= <<<P000979><
<assignableExpression> <assignmentOperator> <expressionWithoutCascade>
<<P000980> <условное выражение>
<<<P000981> <rowExpressionWithoutCascade>

<expressionList> ::= <expression> (',' <expression>)*

<primary> ::= <Это выражение>
<<P000982> <<<P000983>{} <unconditionalAssignableSelector>
<<P000984> \SUPER{{}<аргумент
<<P000986> <выражение функции>
<<P000987> <буквально>
<<P000988> <идентификатор>
<<P000989> <newExpression>
<<<P000990> <constObjectExpression>
<<<P000991><constructorInvocation>
<<<P000992>> <<выражение>> "

<literal> ::= <nullLiteral>
<<<P000993> <booleanLiteral>
<<<P000994> <нумерация
<<P000995> <stringLiteral>
<<<P000996> <symbolLiter>
<<<P000997> <listLiteral>
<<<P000998> <setOrMapLiteral>
<<P000488>

<<<P000849>%
Выражение <<<P000050> всегда могут быть заключены в скобки,
Но это никогда не оказывает семантического влияния на <<<P000051>.

<<<P000999>{%>>
Однако это может оказать влияние на окружающую экспрессию.
Например, для класса <<<P000606> со статическим методом
<<P000607>, <<P000608> возвращается 42.
<<<P000609> - ошибка времени компиляции.
%
Дело в том, что значение <<<P000610>
определяется в нескольких частях,
Вместо того, чтобы быть строго композиционным.
Конкретно, значение <<<P000611> и \code{(C)} Как выражения одинаковы,
но значение <<<P000613> не определяется в терминахЗначение <<<P000614> как выражение,
и отличается от значения <<<P000615>.%
• строго композиционная оценка всегда оценивала бы каждый
% субэкспрессия с использованием одних и тех же правил («оценка» всегда одно и то же);
%, а затем объединит результаты оценки в результат
% от общего выражения. Мы не будем расширять это здесь, и в частности мы не будем
% обсудить, имеет ли композиционная оценка смысл в контексте
% побочных эффектов. Но все же полезно помнить, что у нас есть эти
% «высоко некомпозиционные» элементы в семантике, такие как статические
% метод поиска.
?


<<<P001000> Оценка выражения
<<P000910>

<<<P000850>%
Оценка выражения либо
<<<P001001>{производит объект}{выражение! производит объект
или это
<<<<P001002>{броски}{выражение!броски}
Исключительный объект и связанный с ним след стека.
В первом случае мы также говорим, что выражение
<<<P001003>{оценивает объект}.

<<<P000851>%
При оценке одного выражения <<<P000052>,
определяется в терминах оценки другого выражения <<<P000053>,
обычно подвыражение <<<P000054>,
и оценки <<<P000055> Отбросьте исключение и след стопки,
Оценка <<<P000056> остановиться в этой точке
и бросает тот же самый объект исключения и след стека.


<<<P001004> Идентификация объекта
<<P000911>

<<<P000852>%
Предопределенная функция Дарта <<<P000616>
Определено так, что <<<P000617> Если [f]:

<<P000460>
<<P001005> <<P000059> Оценка нулевого объекта (<<P000519>)
или экземпляра <<<P000618> и <<<P000619>, или
<<P001006> <<<P000062> и <<P000063> являются экземплярами <<<P000620> и <<<P000621>, или
<<P001007> <<P000066> и <<P000067> постоянные строки и <<<P000622>, или
<<P001008> <<<P000070> и <<P000071> Примеры <<<P000623>
и один из следующих пунктов:

<<P000461>
<<P001009> <<P000072> и <<P000073> ненулевой и <<P000624>.
<<<p001010> <<<P000076> и <<P000077> являются <<<P000078>.
<<P001011> Как <<<P000079>, так и <<<P000080> являются <<<P000081>.
<<P001012> <<<P000082> и <<<P000083> представляет собой значение NaN
с тем же базовым битовым рисунком.
<<P000489>
или
<<<P001013> <<P000084> и <<P000085> Постоянные списки, которые определяются как идентичные.
в спецификации буквальных выражений списка (<<<P000520>), ИЛИ
<<<P001014> <<P000086> и <<P000087> Постоянные карты, которые определены как идентичные
в спецификации буквальных выражений карты (<<P000521>), ИЛИ
<<<P001015> <<<P000088> и <<P000089> являются постоянными объектами одного класса <<<P000090>
и значение каждой переменной экземпляра <<<P000091> тождественны
значение соответствующей переменной экземпляра <<<P000092>. или
<<<P001016> <<<P000093> и <<P000094> Это один и тот же объект.
<<P000490>

<<<P001017>{%>
Определение <<<P000625> Двойники отличаются от равенства
в том, что NaN идентичен самому себе,
и что отрицательный и положительный ноль различны.%
?

<<<P001018>{%>
Определение равенства для двойников продиктовано стандартом IEEE 754.
Они утверждают, что NaN не подчиняются закону рефлексивности.
Учитывая, что оборудование реализует эти правила,
Их необходимо поддерживать по соображениям эффективности.

Определение идентичности не ограничено таким же образом.
Вместо этого предполагается, что бит-идентичные двойники идентичны.

Правила идентификации делают невозможным для программиста Дарта
наблюдать, является ли булевое или числовое значение упакованным или неупакованным.%
?


<\subsection{Константы)
<<P000912>

<<<P001020>{%>
Все использования «постоянного» в Дарте связаны со временем компиляции.
Потенциально постоянное выражение — это выражение, которое обычно приводит кпостоянное значение при задании значений определенных параметров.
Постоянные выражения — это подмножество потенциально постоянных выражений.
что \emph{может} оцениваться в момент компиляции.%
?

<<<P001022>{%>
Постоянные выражения ограничиваются выражениями, которые
выполнять только простые арифметические операции, булевы условия;
Творчество и творчество.
Функциональное тело пользователя не выполняется
при постоянной оценке выражения,
Только члены системных классов
<<P000626>, <<P000627>, <<P000628>, <<P000629>,
<<P000630>, <<P000631>, <<<P000632> или <<P000633>.%
?

<<<P000853>%
<<<P001023>{потенциально постоянные выражения}{%
Потенциально постоянное выражение
и \IndexCustom{постоянные выражения}{постоянное выражение}
являются следующие:

<<P000462>
<<<P001025>
Буквальный булев, \TRUE{} или \FALSE{} (<<P000522>),
Это потенциально постоянное и постоянное выражение.

<<P001028>
Буквальное число (<<P000523>) это
Потенциально постоянное и постоянное выражение
если он оценивает тип <<P000634> или <<P000635>.
Слишком большое целое число буквально не оценивает объект.

<<P001029> Буквальная строка (<<P000524>>>) с интерполяцией строк
(<<P000525>)
с выражениями <<P000095>, <<<P001030>, <<p000096> потенциально постоянное выражение
Если <<P000097>, <<<P001031>, <<P000098> Это потенциально постоянные выражения.
Буквальное — это постоянное выражение
Если <<P000099>, <<<P001032>, <<P000100> Постоянные выражения
Оценивая экземпляры <<<P000636>, <<P000637>,
<<P000638>, <<P000639> или <<P000640>.
<<<P001033>{%>
Эти требования носят тривиальный характер
Если в строке нет интерполяций.%
?
<<<P001034>{%>
Было бы заманчиво разрешить интерполяцию струн там, где
интерполированное значение - любая постоянная времени компиляции.
Однако для этого потребуется запустить <<<P000641> метод
для постоянных объектов,
может содержать произвольный код.%
?

<<<P001035>
Буквальный символ (<<P000526>) это
потенциально постоянное и постоянное выражение.

<<<P001036>
Буквальный <<<P001037> (<<P000527>) это
потенциально постоянное и постоянное выражение.

<<P001038>
Идентификатор, обозначающий постоянную переменную, является
потенциально постоянное и постоянное выражение.

<<P001039>
Квалифицированная ссылка на статическую постоянную переменную
(<<P000528>)
который не соответствует отложенному префиксу,
Это потенциально постоянное и постоянное выражение.
<<<P001040>{%>
Например, если класс <<<P000101> объявляет постоянную статическую переменную <<<P000102>,
<<P000642> Это постоянная.
То же самое верно, если <<P000105> Доступ осуществляется через префикс <<<P000106>;
<<P000643> является постоянной, если <<<P000110> Отложенный префикс.%
?

<<P001041>
простой или квалифицированный идентификатор, обозначающий класс;
микширование или псевдоним типа, который не квалифицирован отложенным префиксом,
Это потенциально постоянное и постоянное выражение.
<<<P001042>{%>
Постоянное выражение всегда оценивается как <<<P000644> объект.
Например, если <<<P000111> Это имя класса или типа псевдонима,
Выражение <<P000645> является постоянным,
Если <<P000113> импортируется с приставкой <<<P000114>, <<P000646> это
Константа <<P000647> Пример, представляющий тип <<<P000117>
Если только <<P000118> Отложенный префикс.%
?

<<P001043>
<<<P000119> быть простым или квалифицированным идентификатором
Функция верхнего уровня (<<P000529>)
Статический метод (<<P000530>)
Это не квалифицируется отложенным префиксом.
Если <<P000120> не подвержена генерической функции
(<<P000531>)
P000121 является потенциально постоянным и постоянным выражением.
Если инстанциация общей функции применяется к <<<P000122>и приведенные фактические аргументы типа <<<P001044> T}{1}{s}
<<<P000123> потенциально постоянное и постоянное выражение
если {}f каждый <<<P000124>, является выражением постоянного типа
(<<P000532>).

<<P001045>
Выражение идентификатора, обозначающее параметр константного конструктора
(<<P000533>)
что происходит в списке инициализаторов конструктора,
Это потенциально постоянное выражение.

<<P001046>
Постоянное выражение объекта (<<P000534>)
Это потенциально постоянное и постоянное выражение.

<<P001047>
Постоянный список буквальный (<<P000535>)
Это потенциально постоянное и постоянное выражение.

<<<P001048>
Постоянный набор буквальный (<<P000536>)
Это потенциально постоянное и постоянное выражение.

<<P001049>
Постоянная карта (<<P000537>)
Это потенциально постоянное и постоянное выражение.

<<P001050>
Выражение скобки <<<P000648> это
Потенциально постоянное выражение
Если <<P000126> Это потенциально постоянное выражение.
Это также постоянное выражение, если <<<P000127> Это постоянное выражение.

<<P001051>
Выражение формы <<<P000649> это
Потенциально постоянное выражение
Если <<<P000130> и <<<P000131> Потенциально постоянные выражения
<<<P000650>> Статически связан с
Предопределенная функция dart <<<P000651> обсуждавшееся выше
(P000538).
Кроме того, это постоянное выражение
Если <<<P000132> и <<<P000133> Это постоянные выражения.

<<P001052>
Выражение формы <<<P000652> это
эквивалентно <<P000653> во всех отношениях,
В том числе, является ли она потенциально постоянной или постоянной.

<<P001053>
Выражение формы <<<P000654> потенциально постоянный
Если <<<P000140> и <<<P000141> Оба эти выражения потенциально являются постоянными.
Он также постоянен, если оба <<<P000142> и <<<P000143> являются постоянными и
<<<P000144> Оценивается на пример <<<P000655>
Или пример, который имеет примитивное равенство.
(<<P000539>),
<<<P000145> Оценка нулевого объекта
(<<P000540>).

<<<P001054>
Выражение формы <<<P000656> потенциально постоянный
Если <<P000147> Потенциально постоянный.
Постоянна, если <<<P000148> Это постоянное выражение, которое оценивает
Пример типа <<P000657>.

<<P001055>
Выражение формы <<<P000658> это
Потенциально постоянный, если <<<P000151> и <<<P000152>
Оба эти выражения потенциально являются постоянными.
Постоянна, если <<P000153> является постоянным выражением и
<<P000463>
<<<P001056><<P000154> Для сравнения: <<<P001057>, или
<<<P001058> <<P000155> \TRUE{} и <<P000156>>> Это постоянное выражение
Оценивается на примере типа <<P000659>.
<<P000491>

<<<P001060>
Выражение формы <<<P000660> это
Потенциально постоянный, если <<<P000159> и <<<P000160>
Оба эти выражения потенциально являются постоянными.
Постоянна, если <<<P000161> является постоянным выражением и
<<P000464>
<<<P001061><<P000162> (b) оценивает до <<<P001062>, или
<<P001063> <<P000163> \FALSE{} и <<P000164>>> Это постоянное выражение
Оценивается на примере типа <<P000661>.
<<P000492>

<<P001065>
Выражение формы <<<P000662> это
Потенциально постоянное выражение
Если <<P000166> Это потенциально постоянное выражение.
Это также постоянное выражение, если <<<P000167> это
Постоянное выражение, оценивающее пример типа <<P000663>
Таким образом, <<<P001066>{} обозначает вызов оператора экземпляра.

<<P001067> выражение одной из форм <<<P000664>,
<<P000665> или <<P000666> потенциально постоянный
<<<P000174> и <<<P000175> Оба эти выражения потенциально являются постоянными.
Он также постоянен, если оба <<<P000176> и <<<P000177> Это постоянные выражения, которыеоба оцениваются по экземплярам <<<P000667>,
или оба экземпляра <<<P000668>,
Чтобы оператор символизировал
<<<P001068><<<P001482>> \lit{ |}, соответственно <<P001070>>>{<<P001483>>>}
обозначает вызов оператора экземпляра.

<<<P001071> Выражение одной из форм
<<P000669>, <<P000670>,
<<<P000671> потенциально постоянный
Если <<<P000184> и <<<P000185> Оба эти выражения потенциально являются постоянными.
Он также постоянен, если оба <<<P000186> и $e_2$ Это постоянные выражения, которые
оба оценивают на пример <<<P000672>,
Чтобы оператор символизировал
<<<P001072><<<P001073>> \lit{<<P001075>>>}, соответственно <<P001076>>>{<<P001077>>>}
обозначает вызов оператора экземпляра.

<<P001078> Выражение формы <<<P000673> это
Потенциально постоянное выражение, если <<<P000190> и <<<P000191>
Оба эти выражения потенциально являются постоянными.
Кроме того, это постоянное выражение
<<<P000192> и <<<P000193> Постоянные выражения
и либо оба оценивают на пример <<<P000674> или <<<P000675>,
или оба оценивают на пример <<<P000676>,
<<<P001079>{+} обозначает вызов оператора экземпляра.

<<P001080>
Выражение формы <<<P000677> потенциально постоянное выражение
Если <<P000195> Это потенциально постоянное выражение.
Это также постоянное выражение, если <<<P000196> Это постоянное выражение, которое
Оценка для экземпляра типа <<P000678> или <<P000679>,
Таким образом, <<<P001081>{-} обозначает вызов оператора экземпляра.

<<P001082> Выражение формы <<<P000680>, <<P000681>,
<<P000682>, <<P000683>, <<P000684>,
<<P000685>, <<P000686>, <<P000687>, или
<<P000688>
потенциально постоянный
Если <<<P000215> и <<<P000216> Оба эти выражения потенциально являются постоянными.
Он также постоянен, если оба $e_1$ и $e_2$ Это постоянные выражения, которые
оценивать по экземплярам <<<P000689> или <<<P000690>,
что данный символ оператора обозначает
Воззвание оператора экземпляра.

<<P001083> Выражение формы <<<P000691>
Потенциально постоянная, если <<<P000222>, <<P000223> и <<P000224>>>
Все это потенциально постоянные выражения.
Постоянна, если <<<P000225> является постоянным выражением и
<<P000465>
<<P001084> <<P000226> <<<P001085>{} и <<P000227> является постоянным выражением, или
<<P001086> <<P000228> <<<P001087>{} и <<P000229> является постоянным выражением.
<<P000493>

<<P001088> Выражение формы <<<P000692> потенциально постоянный
Если <<<P000232> и <<<P000233> Оба эти выражения потенциально являются постоянными.
Он постоянен, если <<<P000234> является постоянным выражением и
<<P000466>
<<<P001089> <<P000235> оценивает объект, который не является нулевым объектом, или
<<<P001090> <<P000236> оценивает нулевой объект и <<<P000237> Это постоянное выражение.
<<P000494>

<<P001091> Выражение формы <<<P000693> потенциально постоянный
Если <<P000239> Это потенциально постоянное выражение.
Постоянна, если <<<P000240> Это постоянное выражение, которое
Для примера <<P000694>,
<<<P000695> Обозначает вызов Геттера.

<<P001092> Выражение формы <<<P000696> потенциально постоянный
Если <<P000242> Это потенциально постоянное выражение.
Постоянна, если <<<P000243> Это постоянное выражение, которое
оценивает до нуля, или оценивает до экземпляра <<<P000697>
<<P000698> Обозначает вызов Геттера.

<<P001093> Выражение формы <<<P000699> потенциально постоянный
Если <<P000246> потенциально постоянное выражение
<<P000247> Постоянное выражение типа,и она является постоянной, если <<<P000248> Постоянная.
<<<P001094>{%>
Ошибка времени компиляции для оценки постоянного выражения
Если бы броська была, то есть
Если <<P000249> Оценивает объект, который не является нулевым объектом
и не по типу <<<P000250>>%
?

<<P001095> Выражение формы <<<P000700> потенциально постоянный
Если <<P000253> потенциально постоянное выражение
<<P000254> Постоянное выражение типа,
и она является постоянной, если <<<P000255> Постоянная.

<<P001096>
Выражение формы <<<P000701>
эквивалентен <<<P000702> во всех отношениях,
В том числе, является ли она потенциально постоянной или постоянной.
<<P000495>

<<<P000854>%
А.
<<P000571>
является одним из:

<<P000467>
<<P001097>
Простой или квалифицированный идентификатор
обозначение типовой декларации (типовой псевдоним, класс или смешанная декларация)
который не соответствует отложенному префиксу,
необязательно с последующими аргументами типа формы
<<P000703>
где <<<P000262>, <<<P001098>, <<P000263> Это постоянные выражения типа.
<<P001099>
Тип формы <<<P000704>
где <<<P000265> - выражение постоянного типа.
<<<P001100>
% % TODO(Eernst): Это не позволяет вводить переменные типа
% к самому типу. «Функция <X>(X)» может быть выражением постоянного типа,
%, но это не покрывается действующими правилами: "X" - переменная типа,
% и никогда не допускается.
Тип функции
<<<P001101><<<P000266> Функция <\metavar{typeParameters}>(\metavar{argumentTypes})}
(где $R$ и <<P001104>>>{<\metavar{typeParameters}>} можно опустить)
где <<P000268>, \metavar{typeParameters} и <<P001107>>>{argumentTypes}
(при наличии) содержат только постоянные выражения типа.
<<P001108>
Тип <<<P001109>.
<<P001110>
Тип <<<P001111>.
<<P000496>

% Быть потенциально постоянным — это структурно, а не по типу.
%, но программа по-прежнему должна удовлетворять типизации в сильном режиме.

% Постоянные выражения (например, «const Foo(42)») всегда оцениваются как
% одно и то же значение, причем не более одного значения на местоположение источника.
% Потенциально постоянные выражения, которые не являются постоянными
% допускают простые операции по основным типам (нум, струна, бул, нуль). Они могут
% вычисляется статически без использования пользовательского кода.

Валидно типизированное потенциально постоянное выражение все еще может потерпеть неудачу при оценке.
% Если это происходит при вызывании const, это ошибка времени компиляции.

<<<P000855>%
Ошибка времени компиляции, если требуется выражение
Постоянное выражение,
Но его оценка выльется в исключение.
Ошибка времени компиляции, если утверждение оценивается как часть
постоянная оценка выражения объекта,
И это утверждение станет исключением.

<<<P000856>%
Ошибка времени компиляции, если значение постоянного выражения
зависит от самого себя.

<<<P001112>{%>
В качестве примера рассмотрим: %
?

<<P000002>


<<<P001113> Дополнительные замечания о константах и потенциальных константах
<<P000913>

<<<P001114>{%>
Нет требования, что
Каждое постоянное выражение оценивается правильно.
Когда требуется постоянное выражение
(например, для инициализации постоянной переменной,
или в качестве значения по умолчанию формального параметра,
или в виде метаданных
Мы настаиваем на том, что постоянное выражение
успешно оцениваться во время компиляции.%
?

<<<P001115>{%>
Вышеизложенное не зависит от программы управления потоком.
Простое наличие необходимой постоянной времени компиляции
Оценка, которая провалится в программе, является ошибкой.
Это также рекурсивно:
Поскольку константы соединений состоят из констант,
Если какая-либо часть константы выбрасывает исключение при оценке,
Это ошибка.
С другой стороны, поскольку реализации могут свободно компилировать код с опозданием,
Некоторые ошибки времени компиляции могут проявляться довольно поздно: %
?

<<P000003>

<<<<P001116>{%>
Реализация свободна для немедленного выпускаОшибка компиляции <<<P000705>,
Но это не обязательно делать.
Он может отсрочить ошибки, если не будет немедленно компилироваться.
Деклараций, которые ссылаются на <<P000706>.
Например, он может задержать предоставление ошибки компиляции.
О методе <<<P000707> до первого вызова <<<P000708>.
Однако он не мог выбрать выполнение <<<P000709>,
см. ветвь, которая относится к <<<P000710> не принимается,
2 Успешное возвращение.

Положение в связи с призывом <<<P000711> Это другое.
<<<P000712> не является константой времени компиляции (несмотря на то, что ее значение равно);
При составлении <<<P000713> не нужно давать ошибку времени компиляции.
Реализация может запускать код,
Это приведет к тому, что будет использован геттер для <<<P000714>.
В этот момент инициализация <<P000715> должны иметь место,
который требует компилирования инициализатора,
который вызывает ошибку компиляции.%
?

<<<P001117>{%>
Обработка <<<P000716> заслуживает некоторого обсуждения.
Рассмотрим \code{\NULL{} +2.
Это выражение всегда вызывает ошибку.
Мы могли бы не рассматривать его как постоянное выражение.
(и вообще не допускать <<P000717>) как
Субэкспрессия числовых или булевых постоянных выражений.
Есть два аргумента в пользу его включения:
Во-первых, он постоянен, поэтому мы можем оценить его во время компиляции.
Во-вторых, представляется более полезным дать
ошибка, вытекающая из оценки явно.%
?

<<<P001121>{%>
Можно обоснованно спросить, почему
<<P000269><<P001484> \,\,$e_2$\,:<<P001488>>>$e_3$ и <<P000272>>>\,??? <<<P001490><<<P001491><<<P000273>
имеют постоянные формы.
Если <<P000274> Статически известно, зачем его тестировать?
Ответ заключается в том, что существуют контексты, где <<<P000275> является переменной,
Например, в постоянных инициализаторах конструктора, таких как:
\code{\CONST{} C(foo):\ \THIS.foo = foo??? <<<P001493>someDefault Стоимость: %
?

<<<P001125>{%>
Разница между
потенциально постоянное выражение и постоянное выражение
заслуживает некоторых объяснений.
Ключевой вопрос заключается в том, как относиться к формальным параметрам конструктора.

Если константный конструктор вызывается из постоянного выражения объекта,
Фактические аргументы должны быть постоянными выражениями.
Если бы мы были уверены, что
Константные конструкторы всегда вызывались из постоянных выражений объектов.
Можно предположить, что формальные параметры конструктора
Константы времени компиляции.

Постоянные конструкторы также могут быть использованы для
выражения для создания обычных экземпляров (<<P000541>),
Таким образом, приведенное выше предположение не является в целом обоснованным.

Тем не менее, использование формальных параметров константного конструктора
имеет значительную полезность.
Вводится понятие потенциально постоянных выражений для облегчения
ограниченное использование таких формальных параметров.
В частности, мы разрешаем использование
Формальные параметры константного конструктора
для выражений, в которых участвуют встроенные операторы,
Но не для постоянных объектов, списков и карт.
Например:%
?

<<P000004>

<<<P001126>{%>
Задание на <<<P000718> допускается при допущении
<<<P000719>> постоянный
(перенаправлено с «P000720») Как правило, это не постоянная времени компиляции.
Задание на <<<P000721> аналогично, но вызывает дополнительные вопросы.
В этом случае суперэкспрессия <<P000722> является <<<P000723>,
Для этого необходимо, чтобы <<P000724> быть числовым постоянным выражением
Все выражение должно считаться постоянным.
Формулировка спецификации позволяет предположить
<<<P000725> оценивает в целое число,
Для вызова этого конструктора в постоянном выражении.
Аналогичные аргументы приводятся для <<<P000726> и \code{q}
в задании <<<P000728>.
Однако не допускаются следующие конструкторы: %
?

<<P000005>

<<<P001127>{%>
Проблема в том, что все это противоречит правилам
постоянные списки (<<P000542>),
Карты (<<P000543>),
и объекты (<<P000544>),
все из которых независимо требуют, чтобы их подэкспрессии были
Постоянные выражения.%
?

<<<P001128>{%>
Все нелегальные конструкторы P000729 выше
Не может быть разумно вызван через <<P001129>,
выражение, которое должно быть постоянным, не может зависеть от
Формальный параметр, который может быть или не быть постоянным.
В отличие от этого, юридические примеры имеют смысл независимо от того,
Призывается ли застройщик через \CONST{} или через <<P001131>.

Внимательные читатели, конечно, будут беспокоиться о случаях, когда
Фактическая аргументация в пользу <<P000730> Они являются константами, но не числовыми.
Это исключается правилами о постоянных строителях.
(<<P000545>),
в сочетании с правилами оценки константных объектов (<<P000546>).%
?


<<P001132> Постоянный контекст
<<P000914>

<<<P000857>%
<<<P000276> быть выражением; <<<P000277> происходит в
<<P000572>
если применяется одно из следующих положений:

% Мы избегаем замкнутости "постоянный контекст зависит от постоянного списка буквальных,
% и т.д., что зависит от постоянного контекста, упомянув <<P001133> модификатор
% прямо здесь. Таким образом, «постоянный контекст» является концепцией более низкого уровня.
% на основе синтаксиса и «постоянных выражений X» (например, «постоянный список буквальный»)
% построены на этом.

<<P000468>
<<P001134> <<P000278> является элементом списка или набора буквальных знаков, первый токен которого <<<P001135>,
<<<P000279> является ключом или ценностью входа
Буквальная карта, первым токеном которой является <<P001136>.
<<P001137> <<P000280> Происходит как <<P000731> в конструкцию, полученную из <<<P001138>{metadata}.
<<P001139> <<P000282> является фактическим аргументом в выражении, полученном из
<<<P001140>{constObjectExpression}.
<<P001141> <<P000283> является инициализирующим выражением постоянной переменной декларации
(<<P000547>).
<<P001142> <<P000284> является выражением случая переключения
(P000548).
<<P001143> <<<P000285> - немедленная субэкспрессия
Выражение <<P000286> Это происходит в постоянном контексте.
где <<P000287> это
% % Может быть добавлен позже:
% % не выражение \THROW{} (<<P000549>>>) и
Не функция буквальная
(<<P000550>).
<<P000497>

<<<P001145>{%>
Постоянный контекст вводится в ситуациях, когда
Выражение должно быть постоянным.
Это позволяет использовать <<<P001146> Модификатор должен быть опущен
в случаях, когда он не предоставляет никакой новой информации.%
?


<<<<P001147>{Нулевой}
<<P000915>

<<<P000858>%
Зарезервированное слово <<<P001148> Для сравнения: <<P000573>.

<<P000469>
<nullLiteral> ::= <<<P001149><
<<P000498>

<<<P000859>%
Нулевой объект является единственным экземпляром встроенного класса <<<P000732>.
% Нижеследующее может быть следствием объявления «нулевым»,
%, но мы не разъясняем, мы просто требуем, чтобы это было ошибкой.
Попытка мгновенно <<P000733> вызывает ошибку компиляции времени.
Это ошибка времени компиляции для класса для расширения, смешивания или реализации.
<<P000734>.
<<P000735> Класс расширяется до <<P000736> класс
и не объявляет никаких методов, кроме тех, которые также объявлены <<P000737>.

<<<P001150>{%>
Нулевой объект имеет примитивное равенство.
(<<P000551>)%
?

<<<P000860>%
Статический тип \NULL{} является \code{Null} Тип.


<<<P001152> Номера.
<<P000916>

<<<P000861>%
A <<<P001153>{цифровой буквальный}{буквальный!цифровой}
является либо десятичным, либо шестнадцатеричным числом, представляющим целое число,
или десятичное двойное представление.

<<P000470>
<numericLiteral> ::= <Номер>
<<P001154> <HEX<<<P001494> Номер

<NUMBER> ::= <DIGIT>+ ('.' <DIGIT>+)? <Выражение>?
<<<P001155>> <DIGIT>+ <EXPONENT>

<EXPONENT> ::= ('e' | 'E') ('+' | '-')? <DIGIT>+
<HEX<<<P001495> Номер ::= "0x" <HEX<<<P001496> DIGIT>+
<<P001156> '0X' <HEX<<<P001497> DIGIT>+

<HEX<<<P001498> DIGIT> ::= 'a' ..'f "
<<P001157> "А" .. "
<<P001158> <DIGIT>
<<P000499>

<<<P000862>%
Буквальное число, начинающееся с "0x" или "0X" "
<<<P001159>{hexadecimal integer literal}{literal!hexadecimal integer}.
Имеет числовое целое значение шестнадцатеричного числа
после "0x" (соответственно "0X").

<<<P000863>%
Буквальное число, содержащее только десятичные цифры, является
<<<P001160>{decimal integerliteral}{literal! десятичное целое число.
Он имеет числовое целое значение десятичного числа.

<<<P000864>%
<<<P001161>{integerliteral}{literal!integer}
является либо шестнадцатеричным целым буквальным, либо десятичным целым буквальным.

<<<P000865>%
<<<P000288> быть целым буквальным числом, которое не является операндом
от оператора унарного минуса,
И пусть <<P000289> Статический контекстный тип <<<P000290>.
Если <<P000739> Применяется к <<<P000291> и <<<P000740> не присваивается <<<P000292>,
Статический тип <<<P000293> является <<<P000741>;
Статический тип <<<P000294> является <<<P000742>.

<<<P001162>{%>
Это означает, что целое число буквально обозначает <<<P000743>
когда он удовлетворяет требованиям типа и <<<P000744> Не стал бы.
В противном случае это <<<P000745>, даже в ситуациях, когда это ошибка.
?

<<<P000866>%
Буквальное число, не являющееся целым числом, является
<<<P001163>{double literal}{literal!double}.
<<<P001164>{%>
Двойной буквальный всегда содержит либо десятичную точку, либо экспонентную часть.
?
Статический тип двойного литерала — <<P000746>.

<<<P000867>%
Если <<P000295> является целым буквальным числом с числовым значением <<<P000296> и статический тип <<<P000747>,
<<<P000297> не является операндой унарного оператора минус,
Далее следует оценка <<<P000298> идет следующим образом:

<<P000471>
<<<P001165> Если <<P000299> является шестнадцатеричным целым числом,
<<P000300> <<<P000748>> Класс реализуется как
подписанные 64-битные два целых числа комплемента,
<<<P000301> Оценить на пример <<<P000749> класс
обозначает цифровое значение <<<P000302>,
<<<P001166>< В противном случае <<P000303> Оценивает экземпляр <<<P000750> класс
обозначает цифровое значение <<<P000304>.
Это ошибка времени компиляции, если целое число <<<P000305> не могут быть представлены
Именно на примере <<<P000751>.
<<P000500>

<<<P001167>{%>
Целые числа в Дарте предназначены для реализации как
64-разрядные два дополняют целые представления.
На практике реализация может быть ограничена другими соображениями.
Например, Dart, скомпилированный для JavaScript, может использовать тип номера JavaScript.
эквивалентно Dart <<<P000752>, чтобы представлять целые числа, и если да,
целые буквы с точностью более 53 бит не могут быть представлены
Точно.%
?

<<<P000868>%
Двойная буквальная оценка для экземпляра <<P000753> класс
64-битное число с плавающей точкой двойной точности
В соответствии со стандартом IEEE 754.

<<<P000869>%
Буквальное целое число со статическим типом <<<P000754> и числовое значение <<<P000306>
Оценивается на пример <<<P000755> Класс, представляющий значение <<<P000307>.
Это ошибка времени компиляции, если значение <<P000308>
Не может быть представлен <<<P001168>{точно} экземпляром <<<P000756>.

<<<P001169>{%>
64-битное число с двойной точностью плавающей точки
Обычно используется для представления диапазона реальных чисел.
вокруг точного значения, обозначаемого числом
знак, мантисса и показатель.
Для целых букв, оценивающих до <<P000757>
Значения, которые мы настаиваем на том, что целое число буквального значения
- точное значение <<<P000758> Пример:%
?

<<<P000870>%
Это ошибка времени компиляции для класса для расширения, смешивания или реализации.
<<P000759>.
Это ошибка времени компиляции для класса для расширения, смешивания или реализации.<<P000760>.
Ошибка времени компиляции для любого класса
кроме <<<P000761> и <<P000762> для расширения, смешивания или реализации
<<P000763>.

<<<P000871>%
Примеры <<<P000764> и <<<P000765> все вышеперечисленное
<<<P001170>{==} оператор, унаследованный от <<P000766> класс.


<<<P001171> Булева?
<<P000917>

<<P000872>%
Зарезервированные слова \TRUE{} и <<P001173>>>{} оценивать объекты
<<<P001174>{true}{true, the object} и
<<<P001175>{false}{false, the object}
которые представляют булевы значения истинные и ложные соответственно.
Они являются <<<P001176>{буквальные буквы}{буквальные!булин}.

<<P000472>
<booleanLiteral> ::= <<<P001177><
<<P001178> <<<P001179><
<<P000501>

<<P000873>%
И \NoIndex{true} и <<P001181>>>{false} являются примерами
Встроенный класс <<P000767>,
Нет других объектов, которые реализуют <<P000768>.
Ошибка времени компиляции для класса
расширять, смешивать или реализовывать <<P000769>.

<<<P001182>{%>
\TRUE{} и <<P001184>>>{} Объекты имеют примитивное равенство
(<<P000552>)%
?

<<<P000874>%
Ссылка на геттер <<<P000770> на булево значение возвращается
<<P000771> Объект, который является значением выражения <<<P000772>.
Статический тип булевого литерала — <<P000773>.


<<<P001185> Струны.
<<P000918>

<<P000875>%
A <<<P000574> представляет собой последовательность кодовых единиц UTF-16.

<<<P001186>{%>
Это решение было принято для совместимости с веб-браузерами и Javascript.
Более ранние версии спецификации требовали, чтобы строка была
Последовательность действительных точек кода Unicode.
Программисты не должны зависеть от этого различия.
?

<<P000473>
<stringLiteral> ::= (<multilineString> | <singleLineString>) +
<<P000502>

<<<P000876>%
Струна может быть последовательностью однострочных струн и многострочных струн.

<<P000474>
<singleLineString> ::= <RAW<<<P001499> Песня <<<P001500> LINE<<<P001501> Стрельба
<<P001187> <<<<P001502> LINE<<<P001503> СТРИНГ<<<P001504> SQ<<<P001505> BEGIN<<<P001506> Окончание
<<P001188> <<<<P001507> LINE<<<P001508> СТРИНГ<<<P001509> SQ<<<P001510> BEGIN\_MID> <выражение> \gnewline{}
(<Single<<P001512>) LINE<<<P001513> СТРИНГ<<<P001514> SQ <<<P001515> MID\_MID> <выражение>* <<<P001190><
<<<<P001517> LINE<<<P001518> СТРИНГ<<<P001519> SQ<<<P001520> MID<<<P001521> Окончание
<<P001191> <<<P001522> LINE<<<P001523> СТРИНГ\_ DQ<<<P001525> BEGIN<<<P001526> Окончание
<<P001192> <<<<P001527> LINE<<<P001528> СТРИНГ<<<P001529> DQ<<<P001530> BEGIN\_MID> <выражение <<<P001193><
(<Single<<P001532>) LINE<<<P001533> СТРИНГ <<<P001534> DQ<<<P001535> MID\_MID> <выражение>* \gnewline{}
<<<<P001537> LINE<<<P001538> СТРИНГ<<<P001539> DQ<<<P001540> MID<<<P001541> Окончание

<RAW<<<P001542> Песня <<<P001543> LINE<<<P001544> Строить>::=
<<<P001195>> (<<P001196>) \sq>\<<<P001198>>>>\<<<P001199>>>>)*\sq "
\alt 'r' ''' (<<P001202>>>(''' |<<P001547>>><<P001203>>>'' |\<<<P001204>>>')''

<<<<P001549> Содержание <<<P001550> Коммон: := <<<P001205> \\>>\sq>><<P001208>>>><<P001552>>><<P001207>>>>\<<<P001208>>>>
<<<p001209> <ESCAPE<<<P001554> Последовательность
\alt\\\gtilde('n' | 'r' | 'b' | 't' | 'v' | 'x' | 'u' |<<P001556>>><<P001212>>>><<P001557>>><<P001213>>>)
<<<P001214> <<<<P001558> СТРИНГ<<<P001559> Интерпол

<<<<P001560> СОДЕРЖАНИЕ\_SQ> ::= <STRING\_ СОДЕРЖАНИЕ\_КОММОН> | ''

<<<P001564> LINE<<<P001565> СТРИНГ<<<P001566> SQ<<<P001567> BEGIN\_END::= \gnewline{}
<<<P001216>> <<<<P001569> Содержание <<<P001570> SQ>*<<<P001217> "
<<<P001571> LINE<<<P001572> СТРИНГ<<<P001573> SQ<<<P001574> BEGIN\_MID> ::= \gnewline{}
<<<P001219>> <<<<P001576> СОДЕРЖАНИЕ\_SQ>*$ {''

<<<P001578> LINE<<<P001579> СТРИНГ<<<P001580> SQ<<<P001581> MID\_MID> ::= \gnewline{}
'} <STRING<<<P001583> СОДЕРЖАНИЕ<<<P001584> SQ>*$ {''

<<<<P001585> LINE<<<P001586> СТРИНГ<<<P001587> SQ<<<P001588> MID\_END::= \gnewline{}
'} <STRING<<<P001590> Содержание <<<P001591> SQ>*<<<P001222> "

<<<<P001592> СОДЕРЖАНИЕ\_DQ> ::= <STRING\_ Содержание <<<P001595> КОММОН > | <<<P001223> "

<<<P001596> LINE<<<P001597> Строка <<<P001598> DQ<<<P001599> BEGIN\_END::= \gnewline{}
<STRING<<<P001601> СОДЕРЖАНИЕ\_DQ>* "

<<<<P001603> LINE<<<P001604> СТРИНГ<<<P001605> DQ<<<P001606> BEGIN\_MID ::= \gnewline{}
''' <STRING\_ СОДЕРЖАНИЕ\_DQ>*$ {''

<<<<P001610> LINE<<<P001611> СТРИНГ<<<P001612> DQ<<<P001613> MID\_MID ::= \gnewline{}
'} <STRING\_ СОДЕРЖАНИЕ\_DQ>*$ {''

<<<<P001617> LINE<<<P001618> СТРИНГ<<<P001619> DQ<<<P001620> MID\_END::= \gnewline{}
'} <STRING\_ СОДЕРЖАНИЕ\_DQ>* "
<<P000503>

<<<P000877>%
Одна строка ограничивается
Сопоставление однократных цитат или двойных цитат.

<<<P001228>{%>
Поэтому, <<P000774> и <<P000775> Оба они являются законными.
<<<P000776> и
<<P000777>.
<<<P000778> не является ни действительной строкой, ни <<P000779>.%
?

<<<P001229>{%>
Грамматика гарантирует, что одна строка не может превышать
одна строка исходного кода,
если оно не содержит интерполированного выражения, которое охватывает несколько линий.
?

<<<P000878>%
Соседние струны неявно сцеплены, чтобы сформировать одну строку буквально.

<<<P001230>{%>
Вот пример: %
?

<<P000006>

<<<P001231>{%>
Dart также поддерживает оператор + для конкатенации струн.

Оператор + на струнах требует аргумента струн.
Он не загоняет свой аргумент в нитку.
Это помогает избежать головоломок, таких как %.
?

<<P000007>

<<<P001232>{%>
Что означает «P000780» вместо того, чтобы
<<P000781>.
По соображениям эффективности,
Рекомендуемая идиома Дарта - использование интерполяции строк.%
?

<<P000008>

<<<P001233>{%>
Струнная интерполяция хорошо работает в большинстве случаев.
Главная ситуация, когда она не вполне удовлетворительна
Это для струнных букв, которые слишком велики, чтобы поместиться на линии.
Многолинейные струны могут быть полезны, но в некоторых случаях
Мы хотим визуально выровнять код.
Это можно выразить, написав
меньшие струны, разделенные белым пространством: %
?

<<P000009>

<<<P000879>%
вспомогательный
<<P000575>
поддерживается вне парсера,
Для того, чтобы обеспечить правильное сопоставление струнных интерполяций.

<<<P001234>{%>
Это необходимо, поскольку выражение непростых интерполяций струн
может содержать строковые буквы
с их собственными непростыми интерполяциями.%
?

<<<P000880>%
Правила с именами <<<P001235><<<P001236><<<P001624>BEGIN<<<P001625> Мид.
маркер нажимается на вспомогательный стек, чтобы указать, что
Началась струнная интерполяция данного вида,
где вид \lit{\sq}, \lit{}, \lit{\sqsqsq}, или \lit{"""].
Правила с именами <<<P001243><<<P001244>\_MID<<<P001627> Мид.
Можно использовать только правило с видом сверху вспомогательного стека.
Правила с именами <<<P001245><<<P001246><<<P001628> MID<<<P001629> Конец.
можно использовать только правило с видом сверху вспомогательного стека,
После этого маркер выскакивает.

<<P000475>
<multilineString> ::= <RAW<<<P001630> MULTI<<<P001631> LINE<<<P001632> Стрельба<<P001247> <MULTI<<<P001633> LINE<<<P001634> Строка <<<P001635> SQ<<<P001636> BEGIN<<<P001637> Окончание
<<<P001248> <MULTI<<<P001638> LINE<<<P001639> СТРИНГ<<<P001640> SQ<<<P001641> BEGIN\_MID> <выражение <<<P001249><
(<MULTI<<P001643) LINE<<<P001644> СТРИНГ<<<P001645> SQ<<<P001646> MID\_MID> <выражение>* \gnewline{}
<MULTI<<<P001648> LINE<<<P001649> СТРИНГ<<<P001650> SQ<<<P001651> MID<<<P001652> Окончание
<<P001251> <MULTI<<<P001653> LINE<<<P001654> СТРИНГ<<<P001655> DQ<<<P001656> BEGIN<<<P001657> Окончание
<<P001252> <MULTI<<<P001658> LINE<<<P001659> СТРИНГ<<<P001660> DQ<<<P001661> BEGIN\_MID> <выражение> \gnewline{}
(<MULTI<<P001663>) LINE<<<P001664> СТРИНГ<<<P001665> DQ<<<P001666> MID\_MID> <выражение>* \gnewline{}
<MULTI<<<P001668> LINE<<<P001669> СТРИНГ<<<P001670> DQ<<<P001671> MID<<<P001672> Окончание

<RAW<<<P001673> MULTI<<<P001674> LINE<<<P001675> STRING>::= 'r' '\sqsqsq <<<<P001256> "
\alt 'r' ''''.*? ''''

<QUOTES\_SQ> ::= |\sq >><<<P001259> "

<<<<P001677> СОДЕРЖАНИЕ\_TSQ> ::= \gnewline{}
<QUOTES<<P001679>>>SQ> (<STRING<<P001680>) СОДЕРЖАНИЕ<<<P001681> КОММОН> | ''' | '\<<<P001261>>>" |\<<<P001262>>>"

<MULTI<<<P001684> LINE<<<P001685> СТРИНГ<<<P001686> SQ<<<P001687> BEGIN\_END::= <<<P001263><
<<<P001264>> <<<<P001689> СОДЕРЖАНИЕ<<<P001690> TSQ> <<<P001265> "

<MULTI<<<P001691> LINE<<<P001692> СТРИНГ<<<P001693> SQ<<<P001694> BEGIN\_MID::= <<<P001266><
<<<P001267>> <<<<P001696> СОДЕРЖАНИЕ<<<P001697> TSQ>* <QUOTES\_SQ> '$ {''

<MULTI<<<P001699> LINE\_ СТРИНГ<<<P001701> SQ<<<P001702> MID\_MID ::= \gnewline{}
'} <STRING\_ СОДЕРЖАНИЕ\_TSQ>* <QUOTES\_SQ>$ {''

<MULTI<<<P001707> LINE<<<P001708> СТРИНГ<<<P001709> SQ<<<P001710> MID\_END::= <<<P001269><
'} <STRING\_ СОДЕРЖАНИЕ\_TSQ>*\sqsqsq "

<QUOTES\_DQ> ::= | '''

<<<<P001715> СОДЕРЖАНИЕ\_TDQ> ::= \gnewline{}
<QUOTES\_DQ> (<STRING<<<P001718>) СОДЕРЖАНИЕ<<<P001719> КОММОН> |\sq" |\<<<P001273>>>" |\<<<P001274>>>"

<MULTI<<<P001722> LINE<<<P001723> СТРИНГ\_ DQ<<<P001725> BEGIN\_END::= <<<P001275><
"""" <STRING<<<P001727> СОДЕРЖАНИЕ\_TDQ>*""

<MULTI<<<P001729> LINE<<<P001730> СТРИНГ<<<P001731> DQ<<<P001732> BEGIN\_MID::= <<<P001276><
""" <STRING<<<P001734> СОДЕРЖАНИЕ<<<P001735> TDQ> <QUOTES<<<P001736>DQ>$ {''

<MULTI<<<P001737> LINE<<<P001738> СТРИНГ<<<P001739> DQ<<<P001740> MID\_MID ::= <<<P001277><
'} <STRING<<<P001742> СОДЕРЖАНИЕ<<<P001743> TDQ> <QUOTES<<<P001744>DQ>$ {''

<MULTI<<<P001745> LINE<<<P001746> СТРИНГ<<<P001747> DQ<<<P001748> MID\_END::= \gnewline{}
'} <STRING\_ СОДЕРЖАНИЕ\_TDQ>*""

<ESCAPE<<<P001752> ПОСЛЕДОВАНИЕ > ::=\<<<P001279>>>>>\<\r><<P001755>>><<P001281>>>>\<<<P001282>>>>\<<<P001283>>>>\<<<P001284>>> "
<<<P001285>\<<<P001286>>> <HEX<<<P001760> DIGIT> <HEX<<<P001761> Диджит>
\alt\<<<P001288>>> <HEX<<<P001763> DIGIT > <HEX<<<P001764> DIGIT> <HEX<<<P001765> DIGIT> <HEX<<<P001766> Диджит>
\alt\<<<P001290>>>{) <HEX<<<P001768> DIGIT<<<P001769> Последовательность "

<HEX<<<P001770> DIGIT<<<P001771> Последовательность>::= <<<P001291><
<HEX<<<P001772> DIGIT > <HEX<<<P001773> Цифра? <HEX<<<P001774> Цифра?
\gnewline{} <HEX<<<P001775> Цифра? <HEX<<<P001776> Цифра? <HEX<<<P001777> Цифра?
<<P000504>

<<P000881>%
Многолинейные струны разграничиваются либосовпадающие тройки одиночных котировок или
Сопоставление тройных двойных цитат.
Если первая строка многострочной строки состоит исключительно из
символы белого пространства, определяемые производством <<<p001293>> Белый мир
(<<P000553>),
может быть приставлено \syntax{<<<P001778>>,
И эта линия игнорируется,
включая разрыв линии в конце.

<<<P001295>{%>
Идея состоит в том, чтобы игнорировать первую линию многострочной строки только в белом пространстве.
где белое пространство определяется как вкладки, пространства и окончательный разрыв линии.
Они могут быть представлены непосредственно,
Но так как для большинства персонажей префиксирование с помощью backslash
Идентичность в несырьевой струне,
Мы допускаем и такие формы.%
?

<<P000882>%
В правиле для \synt RAW<<<P001779> MULTI<<<P001780> LINE<<<P001781> Стринг,
Два случая <<<P001297>{.*?}
обозначает шаг нежелательного распознавания токенов:
Он заканчивается, как только Lookahead является указанным следующим токеном.
\commentary{(то есть <<P001299>>>{<<P001300>>>} или <<P001301>>>{"""))}.

<<<P001302>{%>
Обратите внимание, что многострочная интерполяция опирается на
Стек состояния интерполяции вспомогательной струны,
как однострочная интерполяция.%
?

<<P000883>%
Струны поддерживают последовательности побега для специальных персонажей.
Побеги бывают:

<<P000476>
<<P001303>
\syntax{\<<<P001305>>>'} для новой линии, эквивалентной <<P001306>>>{\<<<P001307>>>0A'}.
<<P001308>
\syntax{\<<<P001310>>>'} для возврата перевозки, эквивалентного <<P001311>>>{\<<<P001312>>>0D'}.
<<P001313>
\syntax{\<<<P001315>>>'} для подачи формы, эквивалентной <<P001316>>>{\<<<P001317>>>0C'}.
<<P001318>
\syntax{\<<<P001320>>>'} для заднего пространства, эквивалентно <<P001321>>>{\<<<P001322>>>08'}.
<<P001323>
\syntax{\<<<P001325>>>'} для вкладки, эквивалентной <<P001326>>>{\<<<P001327>>>09'}.
<<P001328>
\syntax{\<<<P001330>>>'} для вертикальной вкладки, эквивалентной <<P001331>>>{\<<<P001332>>>0B'}.
<<P001333>
<<<P001334><<<<P001794><<P001335>>> <HEX<<<P001795> DIGIT><<<P000309> <HEX<<<P001796> DIGIT><<<P000310>, эквивалентный

<<<P001336><<<<P001797><<P001337>< <HEX<<<P001798> DIGIT><<<P000311> <HEX<<<P001799> DIGIT><<<P000312>>
<<P001338>
<<<P001339><<<<P001800><<P001340>> <HEX<<<P001801> DIGIT><<<P000313> <HEX<<<P001802> DIGIT><<<P000314> <HEX<<<P001803> DIGIT><<<P000315> <HEX<<<P001804> DIGIT><<<P000316>,
эквивалентный

<<P001341>
<<<P001342><<<<P001805><<P001343>>><<<P001805> <HEX<<<P001806> DIGIT><<<P000317> <HEX<<<P001807> DIGIT><<<P000318> <HEX<<<P001808> DIGIT><<<P000319> <HEX<<<P001809> DIGIT><<<P000320>>
<<P001344>
<<<P001345><<<<P001810><<P001346>< <HEX<<<P001811> DIGIT<<<P001812> Последовательность > }
Точка кода Unicode, представленная
<<<P001347><<HEX<<<P001813> DIGIT<<<P001814> Последовательность.
Ошибка времени компиляции, если значение
<<<P001348><<HEX<<<P001815> DIGIT<<<P001816> Последовательность
не является действительным кодом Unicode.
<<P001349> Например, }
{\color{комментарийЦвет}{<<<< P001351>>>{'\<<<P001352>>>{0A}'}}\color{нормативный цвет}}
<<<P001354>{является кодовой точкой U+000A. ?
<<P001355>
<<<P001356>\$ Это указывает на начало интерполированного выражения.
<<P001357>
{% Нам нужно определение для <<P000321> Чтобы иметь возможность использовать его в <<<P001358>.
\def\k$k$
В противном случае \syntax{\\\k> указывает характер <<<P001363>< для
любой <<<P001364> не в <<<P001365>{<<P000323>>>'n', 'r', 'f', 'b', 't', 'v', 'x', 'u'<<P000324>>>}.
?
<<P000505>

<<P000884>%
Любая строка может быть приставлена к символу <<<P001366>{r},
Указывая, что это <<<P000576>,В этом случае не распознаются побеги или интерполяции.

<<P000885>%
Линейные разрывы в многострочной строке представлены
<<<P001367> LINE<<<P001820> Производство.
Перерыв линии вводит один новый символ линии (U+000A)
в струнное значение.

<<P000886>%
Это ошибка времени компиляции, если буквальная строка не содержит
символьная последовательность формы \syntax{\<<<P001369>>>'} за которым не следует
последовательность из двух шестидесятых цифр.
Это ошибка времени компиляции, если буквальная строка не содержит
символьная последовательность формы \syntax{\<<<P001371>>>> за которым не следует
либо последовательность из четырех шестидесятых цифр,
или путем кудрявой скобки разграниченной последовательности шестидесятых цифр.

<<P000477>
<LINE<<<P001823> BREAK > ::=\<<<P001372> "
<<<P001373><<<P001825><<P001374><<<P001826><<P001375> "
<<<P001376><<<P001827><<P001377> "
<<P000506>

<<P000887>%
Все строковые буквы оцениваются на экземпляры встроенного класса <<<P000782>.
Ошибка времени компиляции для класса
расширять, смешивать или реализовывать <<P000783>.
<<P000784> класс переопределяет оператор <<<P001378>{==}
<<P000785> класс.
Статический тип строкового литерала — <<<P000786>.


<<<P001379> Струнная интерполяция
<<P000919>

<<<P000888>%
Можно встраивать выражения в несырьевые струнные буквы,
Чтобы эти выражения были оценены,
Полученные объекты преобразуются в струны и
соединяются с ограждающей струной.
Этот процесс называется «P000577».

<<P000478>
<stringInterpolation> ::= <<<<P001828> СТРИНГ<<<P001829> Интерпол
<<<P001380> <${<выражение>> "

<<<<<P001830> СТРИНГ<<<P001831> Интерполация>::= \gnewline{}
"$" (<IDENTIFIER<<P001832>) Нет<<<P001833> DOLLAR > | <BUILT<<<P001834> В <<<P001835> Идентификатор> | \THIS
<<P000507>

<<<P001383>{%>
Читатель заметит, что выражение внутри интерполяции
может включать в себя струны,
которые могут быть интерполированы рекурсивно.%
?

<<<P000889>%
Неудавшийся \lit{\$} символ в строке означает
Начало интерполированного выражения.
<<<P001385><<<P001837>> Знак может сопровождаться либо:

<<P000479>
<<P001386> Единый идентификатор <<<P000835> не содержит \lit{\$} характер
(но это может быть встроенный идентификатор),
или зарезервированным словом <<<P001388>.
<<P001389> Выражение <<P000325> Ограничен кудрявыми брекетами.
<<P000508>

<<<P000890>%
Форма \code{\<<<P000326>>>\{\id\}}.
интерполированная строка <<<P000327>, с содержанием
\code$s_0$\<<<P000329>>>e_1<<P000330>>>>>>_1\ldots{}s_{n-1}<<< P000331>>>\{$e_n$\}$s_{n}$ "
(где любой из <<P000334>) может быть пустой
Оценивается путем оценки каждого выражения <<<P000335> (<<P000336>)
в строку <<<P000337> в порядке, в котором они встречаются в исходном тексте, следующим образом:

<<P000480>
<<P001394> Оценка <<<P000338> Объект <<P000339>.
<<P001395> Нажмите <<<P000787> Метод <<<P000340> без аргументов,
Пусть <<<P000341> быть возвращенным объектом.
<<P001396> Если <<P000342> является нулевым объектом, возникает динамическая ошибка.
<<P000509>

<<<P000891>%
Наконец, результат оценки <<<P000343> это
сцепление струн <<<P000344>, <<P000345>, <<P001397>, <<P000346> и <<P000347>.


<<<P001398>{Символы}
<<P000920>

<<<P000892>%
A <<<P001399>{символ буквальный}{буквальный!символ}
обозначает имя, которое будет
действительное название декларации, действительное название библиотеки или <<<P001400>.

<<P000481>
<symbolLiteral> ::=
'#' (<identifier> ('.' <identifier>)* | <operator> | <<<P001401>)
<<P000510>

<<<P000893>%Статический тип буквального символа <<<P000788>.

<<<P000894>%
<<<P000836> Идентификатор, который не начинается с подчеркивания
(«P000789»).
Буквальный символ <<P000790>
Оценить на пример <<<P000791>
идентификатор <<<P001402>.

<<<P000895>%
Буквальный символ <<<P000792>
где <<<<<<<<<<<<<<<<<<<<<<<<<<<<<<<<<<<<<<<<<<<<<<<<<<<<<<<<<<<<<<<<<<<<<<<<<<<<<<<<<<<<<<<<<<<<<<<<<<<<<<<<<<<<<<<<<<<<<<<<<<<<<<<<<<<<<<<<<<<<<<<<<<<<<<<<<<<<<<<<<<<<<<<<<<<<<<<<<<<<<<<<<<<<<<<<<<<<<<<<<<<<<<<<<<<<<<<<<<<<<<<<<<<<<<<<<<<<<<<<<<<<<<<<<<<<<<<
Для сравнения: <<P000793>
представление этой конкретной последовательности идентификаторов.
<<<P001405>{%>
Этот тип символа буквально обозначает название библиотечной декларации.
Как указано в <<<P001406>{название библиотеки}.
Названия библиотек не подлежат конфиденциальности, даже
если некоторые из его идентификаторов начинаются с подчеркивания.
?

<<<P000896>%
Буквальный символ \code{\#\metavar{op}}
где \metavar{op} происходит от \synt{operator}
Оценить на пример <<<P000794>
Укажите конкретное имя оператора.

<<<P000897>%
Буквальный символ <<<P000795>
Для сравнения: <<P000796>
Представляет собой зарезервированное слово <<<P001411>.

<<<P000898>%
За значение <<P000352> Буквальный символ, представляющий термин исходного кода
Как указано в предыдущих параграфах, мы говорим, что <<P000353> является
<<<P001412>{нечастный символ на основе}{символ!нечастный, на основе}
строка, содержание которой является символами этого термина, без белого пространства.

<<<P001413>{%>
Обратите внимание, что это не относится к частным символам.
которые обсуждаются ниже.
Частный символ не основан на какой-либо строке.
?

<<<P000899>%
Если <<P000354> является значением вызова <<<P000797> конструктор
в форме
<<P000798>, <<P000799>,
или <<P000800>,
где <<P000358> является выражением
\commentary{(постоянно, если необходимо)}
который оценивает строку <<<P000359>,
Мы говорим, что <<P000360> <<<P001415>{неприватный символ на основе} <<P000361>.

<<<<P001416>{%>
Обратите внимание, что <<<P000801> является частным символом,
и отличается от <<<P000802>, как описано ниже.%
?

<<<P000900>%
Предположим, что <<P000362>,
<<<P000363>> значение постоянного выражения
Символ, основанный на строке <<<P000364>.
Если <<<P000803>, то <<<P000367> и <<<P000368> Это один и тот же объект.
<<<P001417> То есть символьные экземпляры канонизированы. ?

<<<P000901>%
Если <<P000369> и <<P000370> Нечастные символы
<<<P001418>{(не обязательно постоянная)}
основанный на струнах <<P000371> и <<P000372>
<<P000373> и <<P000374> равны в соответствии с оператором \lit{===
если и только если <<P000804>
(P000554).

<<<P000902>%
Буквальный символ <<<P000805>, где <<<P000806> Является идентификатором
Оценить на пример <<<P000807>
представляет собой частный идентификатор <<<P000808> библиотеки-приложения.
Все случаи <<<P000809> <<<P001420>{в той же библиотеке} оценить
Тот же объект,
и никаких других символов, буквальных или <<P000810> призыв конструктора
оценивает один и тот же объект,
и не для <<P000811> Пример, равный этому объекту
в соответствии с оператором <<<P001421>{==}.

<<<P001422>{%>
Можно задаться вопросом, какова мотивация введения буквальных символов?
В некоторых языках символы канонизированы, а струны нет.
Однако буквальные струны уже канонизированы в Дарте.
Символы немного легче печатать по сравнению со строками.
Их употребление может стать странным,
Но это не является достаточным основанием для добавления
Буквальная форма языка.
Основная мотивация связана с использованием рефлексии.
Веб-специфическая практика, известная как минификация.

Минификация последовательно сжимает идентификаторы во всей программе.
Чтобы уменьшить размер загрузки.
Эта практика создает трудности для рефлексивных программ.
Это относится к программным декларациям через строки.
строка будет ссылаться на идентификатор в источнике,Но идентификатор больше не будет использоваться в минированном коде.
и отражающий код, использующий их, не сработает.
Поэтому в отражении Дарта используются объекты типа <<<P000812>
Вместо струн.
Случаи <<<P000813> гарантируется стабильная
Что касается минирования.
Предоставляя буквальную форму для символов делает
Отражающий код легче читать и писать.
Символы легко печатаются и часто могут действовать как
Удобные заменители для энумов являются вторичными преимуществами.
?


<\subsection{Сборник литературы}
<<P000921>

<<<P000903>%
В этом разделе указаны несколько буквальных выражений, обозначающих коллекции.
Некоторые синтаксические формы могут обозначать более одного вида коллекции.
В этом случае выполняется этап раздвоения
Для того чтобы определить вид
(P000555).

<<<P000904>%
Подразделы настоящего раздела посвящены механизмам, которые
Общие для всех видов коллекции буквы
(<<P000556>,
<<P000557>,
Далее следует спецификация букварей списка
(<<P000558>, <<P000559>,
Затем следует спецификация того, как различать и выводить типы
для наборов и карт
(<<P000560>, <<P000561>,
И, наконец, спецификация наборов
(<<P000562>)
и карты
(P000563).

<<P000482>
<listLiteral> ::= <<<P001424>> <typeArguments>? <элементы>? "

<setOrMapLiteral> ::= <<<P001425>> <typeArguments>? <элементы>? "

<элементы>::= <элемент> (',' <элемент>)*,'?

<element> ::= <expressionElement>
<<<P001426> <mapElement>
<<<P001427> <spreadElement>
<<P001428> <ifElement>
<<P001429> <для элемента>

"Экспрессии" :: <Экспрессии>

<mapElement> ::= <expression> <expression>

<spreadElement> ::= ('...' | '...?') <выражение>

<ifElement> ::= \IF{} '(' <выражение> ')' <элемент> (<<P001431>{} <элемент>)?

<forElement> ::= <<<P001432> \FOR{} '(' <forLoopParts> ')' <элемент>
<<P000511>

<<<P000905>%
Синтаксически, a <<<P000578> может быть
а \synt{listLiteral} или <<P001435>>>{setOrMapLiteral}.
Содержание коллекции определяется как последовательность
<<<P001436>{сбор буквальных элементов}{сбор буквальных элементов}! элементы},
коротко <<P000579>.
Каждый элемент может быть декларативной спецификацией одного объекта,
Например, \synt{expressionElement} или \synt{mapElement},
может указывать коллекцию, которая должна быть включена,
в форме <<<<P001439>{spreadElement}
или это может быть вычислительный элемент, определяющий
Как получить ноль или больше объектов
в форме \synt{ifElement} или <<P001441>>>{forElement}.

<<<P001442>{%>
Термины, полученные из <<<P001443>{элемента},
Способность создавать коллекции из них;
Также известен как
<<<P000580>.%
?

<<<P000906>%
<<<P000581> элемент <<<P000377> полученный из
\synt{expressionElement} или <<P001445>>>{mapElement}
<<P000378>.
Листовые элементы элемента формы
<<P000814> или
\code\FOR\,(<<P001448>>>{forLoopParts})<<P001847>>>$\ell$}
Это элементы листьев <<<P000382>.
Листовые элементы элемента формы
<<P000815>
Это соединение листовых элементов <<<P000386> и <<<P000387>.
Элементами листа <<<P001449>{spreadElement} является пустой набор.

<<<P001450>{%>
Листовые элементы сборника в буквальном смысле всегда
Набор элементов выражения и/или элементов карты.%
?

<<<P000907>%
Для того, чтобы позволить сбору букв в качестве постоянных выражений,
Уточним, что это означает для элемента <<<P000388>
быть
<<<P001451>{constant}{сборник буквальный элемент!constant}
или
<<<P001452>{потенциально постоянная}{%
Буквальный элемент коллекции! Потенциально постоянные:

<<P000483>
<<P001453>
Когда <<P000389> <<<P001454>{выражениеЭлемент} формы <<P000390>:

<<P000391> потенциально постоянный элемент
Если <<P000392> является потенциально постоянным выражением,
<<P000393> является постоянным элементомЕсли <<P000394> Это постоянное выражение.
<<P001455>
Когда <<P000395> <<<P001456>{mapElement}
в форме <<<P000816>:

<<P000398> потенциально постоянный элемент
<<<P000399> и <<<P000400> потенциально постоянные выражения,
Это постоянный элемент, если они являются постоянными выражениями.
<<P001457>
Когда <<P000401> <<<P001458>{распространенный элемент}
в форме \code{...<<<P000402>>>} или \code{...?<<<P000403>>>}:

<<P000404> потенциально постоянный элемент
Если <<P000405> Это потенциально постоянное выражение.

<<P000406> является постоянным элементом
Если <<P000407> Постоянное выражение, оценивающее
<<<P000819>, <<<P000820> или <<P000821> пример
Первоначально созданный списком, набором или картой буквально.

Кроме того, <<P000408> является постоянным элементом, если это <<<P000822>,
где <<P000410> Постоянное выражение, оценивающее
к нулевому объекту.
<<P001459>
Когда <<P000411> <<<P001460>{ifElement}
в форме
<<P000823>
или форма
<<P000824>:

<<P000417> потенциально постоянный элемент
Если <<P000418> является потенциально постоянным выражением,
<<P000419> потенциально постоянный,
и так далее <<<P000420>, если присутствует.

<<P000421> является постоянным элементом, если <<<P000422> Это постоянное выражение
и:

<<P000484>
<<P001461>
<<<P000423> <<<P000825> и
<<<P000426> <<<P001462>{} и <<P000427> является постоянным,
<<<P000428> <<<P001463>{} и <<P000429> Потенциально постоянный.
<<P001464>
<<<P000430> <<<P000826> и
<<<P000434> Показатель <<<P001465>,
<<P000435> является постоянным,
<<P000436> потенциально постоянный;
<<<P000437> Показатель <<<P001466>,
<<P000438> потенциально постоянный,
<<<P000439> является постоянным.
<<P000512>
<<P000513>

<<<P001467>{%>
A <<<P001468>{forElement} никогда не может происходить в постоянном сборе букв.%
?

% % TODO(Eernst): Мы можем изменить этот текст на комментарий или удалить его,
% состоит из описаний ошибок, но они уже охвачены
% в других местах.
% %
% ошибок компиляции.
% %
% % Следующие ошибки времени компиляции указаны в функции
Спецификация %%, но они уже подразумеваются указанными ошибками
% в других местах.
% %
%% - "Неосведомленный элемент спреда имеет статический тип <<<P000827>".
%% Для списка буквально: ListLiteralInference делает эту ошибку случайностью
Процентный анализ на элементе «Spread».
% % Set/map буквально: setAndMapLiteralInference, ditto.
% %
"Спредный элемент в списке или наборе буквальных элементов имеет статический тип, который является
% % не <<<P001469> и не является подтипом <<<P000828>".
% % Список буквальный: анализ случаев на «элементе спреда».
% % Набор буквальный: анализ кейсов по «элементу спреда»: преуспевает только в том случае, если
% % <<<P000440> реализует «Карту» (а не «Итерируемую»), поэтому имеет ключ/значение
Тип %% и тип элемента без ошибок (P000564).
% %
%% - "Элемент спреда в списке или наборе имеет статический тип, который
% Реализация <<<P000829> <<<P000442> и
%% <<<P000443> не присваивается типу элемента списка".
% % Перечислите буквально: ср. «Элемент распространения» означает, что элемент распространения
% % имеет тип элемента <<<P000444>, а неназначение является ошибкой <<<P000565>.
% % Набор буквальный: cf. 'Spread element' for set/maps: ditto.
% %
%% - "Спредный элемент на карте в буквальном смысле имеет статический тип, который является
% % не <<<P001470> и не является подтипом <<<P000830>".
% % Анализ кейсов «Spread element» подразумевает, что он реализует
% % «Изменяемый», а не «Карта», в этом случае у него нет пары типа ключ/значение,
%, что является ошибкой, <<P000566>.
% %
%% - "Если статичный тип элемента карты распространяется <<<P000831>
% для некоторых <<<P000447> и <<<P000448> и <<<P000449> не присваивается ключевому типуКарта %% или <<P000450> не присваивается типу значения карты".
% % Это ошибка, <<P000567>.
% %
%% - "Переменная в <<<P001471>{для элемента} <<<P001472>{(либо для входа)
%% элемент или C-стиль)} объявляется вне элемента окончательным или
% % не иметь сеттера».
% % Эта ошибка указана в <<<P000568> и
%% <<<P000569>.
% %
%% - "Тип выражения итератора в синхронном элементе включения
% не может быть присвоено <<<P000832> Для любого типа <<P000452>.
В противном случае <<<P000582> итератора составляет <<<P000453>.
%, покрытых тем же текстом, что и предыдущая ошибка.
% %
%% - "Переходный тип итератора в синхронном элементе включения
% не может быть отнесено к типу переменной "for-in".
%, покрываемых одним и тем же текстом.
% %
%% - "Тип выражения потока в асинхронном \AWAIT{} вход
Элемент %% не может быть присвоен <<<P000833> Для любого типа <<P000455>.
% % В противном случае <<<P000583> Поток имеет значение <<<P000456>".
%% Один и тот же текст (текст «для элемента» включает в себя как ожидание, так и ожидание).
% %
%% - "Тип потока итератора в асинхронном \AWAIT{} элемент for-in
% не может быть отнесено к типу переменной "for-in".
% опять тот же текст.
% %
%% - "\AWAIT{} <<<<P001476>< Применяется, когда функция непосредственно заключает
Буквальный набор %% не является асинхронным".
% опять тот же текст.
% %
% % — «\AWAIT Используется перед C-стилем <<<P001478>{для элемента}.
%% <<<P001479> Его можно использовать только в петлях ".
% опять тот же текст.
% %
%% - "Тип выражения состояния \commentary{(второй пункт)}
% в стиле C \FOR{} Элемент не может быть присвоен <<<P000834>".
% опять тот же текст.


